%% toc: P15 157 157.1,157.2,157.3,157.4
% Proofs for the explicit moments, L^p inequalities and integral approximations (Durrett, secs. 1.4-1.6).
\section{Explicit moments, norm inequalities and integral approximations: proofs}
\label{sec:annex-p15-w02b-p15-001-poisson-geometric}

Use the definitions of
Chapter~\ref{sec:p15-w02b-p15-001-poisson-geometric}. Two facts about
nonnegative sums are used repeatedly. First, by
\cref{def:p15-measure}, a nonnegative sum is the supremum of its
finite partial sums, so relabelling the index set by a bijection does
not change it, and terms equal to zero may be added or removed.
Second, for a finite \(c\geq0\),
\(\sum_ic\,a_i=c\sum_ia_i\): each finite sum satisfies it, for
\(c>0\) multiplication by \(c\) is an order-preserving bijection of
\([0,\infty]\) and so preserves suprema, and for \(c=0\) both sides
are zero by the convention \(0\cdot\infty=0\).

\begin{annexitem}{annex:p15-w02b-p15-001-uniform-moments}{thm:p15-w02b-p15-001-uniform-moments}{Sums of integers and squares by induction, and the uniform moments}
\begin{proof}
(1) It suffices to prove, for every integer \(m\geq0\),
\[
 \sum_{k=1}^{m}k=\frac{m(m+1)}2,\qquad
 \sum_{k=1}^{m}k^2=\frac{m(m+1)(2m+1)}6,
\]
where a sum over an empty range is zero. Both identities hold at
\(m=0\), where every side is zero. Suppose they hold at \(m\). Then
\[
 \sum_{k=1}^{m+1}k=\frac{m(m+1)}2+(m+1)=\frac{(m+1)(m+2)}2,
\]
and
\[
 \begin{aligned}
 \sum_{k=1}^{m+1}k^2&=\frac{m(m+1)(2m+1)}6+(m+1)^2\\
 &=\frac{(m+1)\bigl(m(2m+1)+6(m+1)\bigr)}6\\
 &=\frac{(m+1)(2m^2+7m+6)}6=\frac{(m+1)(m+2)(2m+3)}6,
 \end{aligned}
\]
since \((m+2)(2m+3)=2m^2+7m+6\). The right sides are the two formulas
at \(m+1\), because \(2(m+1)+1=2m+3\). By induction both identities
hold for every \(m\geq0\); in particular the two right-hand identities
of part~(1) hold for \(m\geq1\).

For the sums from zero, let \(m\geq1\). The term with \(k=0\) is zero,
so \(\sum_{k=0}^{m-1}k=\sum_{k=1}^{m-1}k\) and
\(\sum_{k=0}^{m-1}k^2=\sum_{k=1}^{m-1}k^2\). Apply the identities just
proved at \(m-1\geq0\):
\[
 \begin{aligned}
 \sum_{k=1}^{m-1}k&=\frac{(m-1)m}2,\\
 \sum_{k=1}^{m-1}k^2&=\frac{(m-1)m\bigl(2(m-1)+1\bigr)}6
 =\frac{(m-1)m(2m-1)}6.
 \end{aligned}
\]

(2) Fix \(m\geq1\). For \(0\leq k\leq m-2\),
\(C_k=(k/m,(k+1)/m]\), since this interval lies in \((0,1)\); and
\(C_{m-1}=((m-1)/m,1)\). Each is a Borel subset of \((0,1)\), hence an
event of \(\mathcal F_U\), and each is a bounded interval of length
\(1/m\), so \(\mathbb U(C_k)=1/m\) by
\cref{thm:p15-lebesgue-unit-interval}, which gives every bounded interval
its length whichever endpoints are included. The intervals
\((k/m,(k+1)/m]\) for distinct \(k\) are disjoint, so the cells are
pairwise disjoint. For \(u\in(0,1)\) let \(k\) be the least integer with
\(mu\leq k+1\); then \(k\geq0\) because \(mu>0\), \(k\leq m-1\) because
\(mu<m\), and \(k<mu\) by the minimality of \(k\). Thus
\(k/m<u\leq(k+1)/m\) and \(u\in C_k\). The union of the cells is
\((0,1)\).

On \(C_k\) the three functions \(s_m,U,t_m\) take the values
\(k/m\), \(u\) and \((k+1)/m\), with \(0\leq k/m<u\leq(k+1)/m\). Hence
\(0\leq s_m\leq U\leq t_m\), and squaring nonnegative numbers preserves
order, so \(s_m^2\leq U^2\leq t_m^2\). Because the cells partition
\((0,1)\), \(s_m^2=\sum_k(k/m)^2\mathbf1_{C_k}\) and
\(t_m^2=\sum_k((k+1)/m)^2\mathbf1_{C_k}\) are nonnegative simple
functions in the sense of \cref{def:p15-simple-integral}. By that
definition and the identities of part~(1),
\[
 \begin{aligned}
 \int s_m\,d\mathbb U&=\sum_{k=0}^{m-1}\frac km\cdot\frac1m
   =\frac{(m-1)m}{2m^2}=\frac{m-1}{2m},\\
 \int t_m\,d\mathbb U&=\frac1{m^2}\sum_{j=1}^{m}j
   =\frac{m(m+1)}{2m^2}=\frac{m+1}{2m},
 \end{aligned}
\]
where the second sum is reindexed by \(j=k+1\). Likewise
\[
 \begin{aligned}
 \int s_m^2\,d\mathbb U&=\frac1{m^3}\sum_{k=0}^{m-1}k^2
   =\frac{(m-1)m(2m-1)}{6m^3},\\
 \int t_m^2\,d\mathbb U&=\frac1{m^3}\sum_{j=1}^{m}j^2
   =\frac{m(m+1)(2m+1)}{6m^3}.
 \end{aligned}
\]
The variables \(U\) and \(U^2\) are nonnegative, so their expectations
are their nonnegative integrals (\cref{def:p15-expectation-integral}),
and the monotonicity of the nonnegative integral
(\cref{thm:p15-monotone-integral}) gives the four displayed bounds.

(3) Expanding the bounds,
\[
 \frac{m-1}{2m}=\frac12-\frac1{2m},\qquad
 \frac{m+1}{2m}=\frac12+\frac1{2m},
\]
so \(|\mathbb E[U]-1/2|\leq1/(2m)\) for every \(m\geq1\). The left
side does not depend on \(m\), and \(1/(2m)\to0\); hence
\(\mathbb E[U]=1/2\). Similarly
\((m-1)m(2m-1)=2m^3-3m^2+m\) and \(m(m+1)(2m+1)=2m^3+3m^2+m\), so
\[
 \begin{aligned}
 \frac{(m-1)m(2m-1)}{6m^3}&=\frac13-\frac1{2m}+\frac1{6m^2},\\
 \frac{m(m+1)(2m+1)}{6m^3}&=\frac13+\frac1{2m}+\frac1{6m^2}.
 \end{aligned}
\]
Both bounds lie within \(1/(2m)+1/(6m^2)\leq2/(3m)\) of \(1/3\), so
\(|\mathbb E[U^2]-1/3|\leq2/(3m)\) for every \(m\geq1\), and
\(\mathbb E[U^2]=1/3\). In particular \(\mathbb E[U^2]<\infty\), and
\cref{thm:p15-expectation-bounds} gives
\(\operatorname{Var}(U)=\mathbb E[U^2]-(\mathbb E[U])^2
 =1/3-1/4=1/12\).

For the remark after the theorem, the expansions above show that in
both cases the upper bound exceeds the lower one by
\(1/(2m)+1/(2m)=1/m\). At
\(m=10\) the bounds for the second moment are
\(9\cdot10\cdot19/6000=0.285\) and \(10\cdot11\cdot21/6000=0.385\).
\end{proof}
\end{annexitem}

\begin{annexitem}{annex:p15-w02b-p15-001-series-tools}{thm:p15-w02b-p15-001-series-tools}{Expectations as series and the negative-binomial series}
\begin{proof}
(1) Singletons and the countable set \(S\) are Borel, so every event
\(\{X=s\}\) and \(\{X\in S\}\) is measurable, and \(h(X)\) is
measurable by \cref{thm:p15-measurable-operations}. Enumerate \(S\)
as \(s_1,s_2,\ldots\) (a finite list if \(S\) is finite). The events
\(\{X=s_n\}\) are pairwise disjoint with union \(\{X\in S\}\). Hence,
at every outcome,
\[
 h(X)\mathbf1_{\{X\in S\}}=\sum_{n}h(s_n)\mathbf1_{\{X=s_n\}}:
\]
at an outcome with \(X=s_m\) only the \(m\)-th term is nonzero, and
off \(\{X\in S\}\) every term is zero. Each term is a nonnegative
simple function with integral \(h(s_n)\mathbb P(X=s_n)\). Termwise
integration of a nonnegative series
(\cref{thm:p15-nonnegative-integral-laws}) gives
\[
 \mathbb E\bigl[h(X)\mathbf1_{\{X\in S\}}\bigr]
   =\sum_nh(s_n)\mathbb P(X=s_n)=\sum_{s\in S}h(s)\mathbb P(X=s),
\]
the last equality by relabelling. Since
\(\mathbb P(X\notin S)=0\), the functions \(h(X)\) and
\(h(X)\mathbf1_{\{X\in S\}}\) agree a.s., so their integrals agree by
part~(2) of the same theorem. This proves the first identity.

Now let \(h\) be real-valued Borel with \(h\geq0\) on \(S\). Its
negative part \(h^-\) vanishes on \(S\), so
\(h(X)^-=h^-(X)\) vanishes on \(\{X\in S\}\) and is zero a.s.;
its expectation is zero. The expectation of \(h(X)\) is therefore
defined and equals \(\mathbb E[h^+(X)]\). The function \(h^+\) is
nonnegative everywhere and equals \(h\) on \(S\), so the first part
gives \(\sum_{s\in S}h^+(s)\mathbb P(X=s)=\sum_{s\in S}h(s)\mathbb P(X=s)\).

(2) The case \(r=0\). For every integer \(M\geq0\), the finite sum
\(T_M=\sum_{m=0}^Mq^m\) satisfies \((1-q)T_M=1-q^{M+1}\), since the
intermediate powers cancel. Thus \(T_M=(1-q^{M+1})/p\). If \(q=0\),
then \(q^{M+1}=0\). If \(0<q<1\), write \(q=1/(1+\delta)\) with
\(\delta>0\). Induction gives Bernoulli's inequality
\((1+\delta)^n\geq1+n\delta\): it holds at \(n=0\), and
\((1+\delta)^{n+1}\geq(1+n\delta)(1+\delta)\geq1+(n+1)\delta\).
Hence \(0\leq q^n\leq1/(1+n\delta)\to0\). The nonnegative sum is the
increasing limit of its partial sums, so
\(\sum_{m\geq0}q^m=\lim_MT_M=1/p\).

A summation identity. For integers \(r\geq0\) and \(m\geq0\),
\[
 \sum_{i=0}^m\binom{i+r}{r}=\binom{m+r+1}{r+1}.
\]
First, Pascal's rule: for integers \(0\leq j\leq n-1\),
\[
 \begin{aligned}
 \binom nj+\binom n{j+1}
 &=\frac{n!}{j!\,(n-j)!}+\frac{n!}{(j+1)!\,(n-j-1)!}\\
 &=\frac{n!\,\bigl((j+1)+(n-j)\bigr)}{(j+1)!\,(n-j)!}
 =\binom{n+1}{j+1}.
 \end{aligned}
\]
The summation identity holds at \(m=0\), where both sides equal one.
If it holds at \(m\), then the sum to \(m+1\) is
\(\binom{m+r+1}{r+1}+\binom{m+r+1}{r}\), which is
\(\binom{m+r+2}{r+1}\) by Pascal's rule with \(n=m+r+1\) and
\(j=r\leq n-1\).

The induction on \(r\). Suppose the identity of part~(2) holds at
\(r\). Consider the countable index set
\(I=\{(i,m):0\leq i\leq m\}\) with nonnegative terms
\(a_{(i,m)}=\binom{i+r}{r}q^m\). Partition \(I\) by the value of
\(m\). The block of a fixed \(m\) is finite, and by the summation
identity its sum is \(\binom{m+r+1}{r+1}q^m\). Partition \(I\)
instead by the value of \(i\). The block of a fixed \(i\) has sum
\(\binom{i+r}{r}\sum_{m\geq i}q^m\). The relabelling \(m=i+n\) and
the scaling remark give \(\sum_{m\geq i}q^m=q^i\sum_{n\geq0}q^n=q^i/p\).
Both partitions are countable families of pairwise disjoint sets with
union \(I\), so \cref{thm:p15-countable-atomic} equates each grouped
sum with \(\sum_{(i,m)\in I}a_{(i,m)}\). Therefore
\[
 \sum_{m\geq0}\binom{m+r+1}{r+1}q^m
 =\frac1p\sum_{i\geq0}\binom{i+r}{r}q^i
 =\frac1p\,p^{-(r+1)}=p^{-(r+2)},
\]
using the induction hypothesis. No convergence hypothesis is needed at
any step, because every sum is a sum of nonnegative terms.
\end{proof}
\end{annexitem}

\begin{annexitem}{annex:p15-w02b-p15-001-poisson-moments}{thm:p15-w02b-p15-001-poisson-moments}{Poisson factorial moments by an index shift}
\begin{proof}
Let \(\lambda>0\), \(S=\{0,1,2,\ldots\}\) and, for a fixed integer
\(k\geq1\), \(h_k(x)=x(x-1)\cdots(x-k+1)\). This polynomial is
continuous, hence Borel, and it is nonnegative on \(S\)
(\cref{def:p15-w02b-p15-001-laws}). Since
\(\mathbb P(X\in S)=1\), the variable \(h_k(X)\) is nonnegative a.s.,
and part~(1) of \cref{thm:p15-w02b-p15-001-series-tools} gives
\[
 \mathbb E\bigl(X(X-1)\cdots(X-k+1)\bigr)
 =\sum_{j\geq0}j(j-1)\cdots(j-k+1)\,e^{-\lambda}\frac{\lambda^j}{j!}.
\]
The terms with \(j<k\) are zero and are removed. For \(j\geq k\) the
descending product consists of the \(k\) largest factors of \(j!\), so
it equals \(j!/(j-k)!\), and the \(j\)-th term is
\(\lambda^k\cdot e^{-\lambda}\lambda^{j-k}/(j-k)!\). Relabel by
\(i=j-k\) and apply the scaling remark:
\[
 \mathbb E\bigl(X(X-1)\cdots(X-k+1)\bigr)
 =\lambda^k\sum_{i\geq0}e^{-\lambda}\frac{\lambda^i}{i!}=\lambda^k.
\]
The last sum is the total Poisson mass, which is one by
\cref{thm:p15-stirling-poisson-premises}.

With \(k=1\) and \(k=2\), \(\mathbb EX=\lambda\) and
\(\mathbb E[X(X-1)]=\lambda^2\). On \(S\), \(x^2=x(x-1)+x\) with
both summands nonnegative. Outside the null event \(\{X\notin S\}\),
\(X^2=X(X-1)\mathbf1_{\{X\in S\}}+X\mathbf1_{\{X\in S\}}\) with both
summands nonnegative. Hence the a.s.\ equality of integrals and
nonnegative additivity (\cref{thm:p15-nonnegative-integral-laws})
give \(\mathbb EX^2=\lambda^2+\lambda<\infty\), where the two
restricted expectations equal the unrestricted ones by the same
a.s.\ equality. Since \(\mathbb EX^2\) is finite,
\cref{thm:p15-expectation-bounds} gives
\(\operatorname{Var}(X)=\mathbb EX^2-(\mathbb EX)^2
 =\lambda^2+\lambda-\lambda^2=\lambda\).

For the point mass at zero, \(X=0\) a.s., so each descending product
of order \(k\geq1\) is zero a.s. (its factor \(X\) is zero), its
expectation is \(0=0^k\), and \(\operatorname{Var}(X)=0\).
\end{proof}
\end{annexitem}

\begin{annexitem}{annex:p15-w02b-p15-001-geometric-moments}{thm:p15-w02b-p15-001-geometric-moments}{Geometric moments without term-by-term differentiation}
\begin{proof}
Let \(0<p<1\) and \(q=1-p\), so \(0<q<1\).

(1) Relabel by \(m=k-1\) and scale:
\(\sum_{k\geq1}pq^{k-1}=p\sum_{m\geq0}q^m=p\cdot(1/p)=1\), by
part~(2) of \cref{thm:p15-w02b-p15-001-series-tools} at \(r=0\).

(2) Fix an integer \(r\geq1\) and let \(S=\{1,2,\ldots\}\) and
\(h_r(x)=x(x-1)\cdots(x-r+1)\), a Borel function nonnegative on
\(S\). By part~(1) of \cref{thm:p15-w02b-p15-001-series-tools},
\(N(N-1)\cdots(N-r+1)\) is nonnegative a.s. and
\[
 \mathbb E\bigl(N(N-1)\cdots(N-r+1)\bigr)
  =\sum_{k\geq1}k(k-1)\cdots(k-r+1)\,pq^{k-1}.
\]
If \(1\leq k\leq r-1\), the factor \(k-i\) with \(i=k\) is zero, so
these terms are removed. If \(k\geq r\), the descending product equals
\(k!/(k-r)!=r!\binom kr\). Relabel by \(k=m+r\) with \(m\geq0\); since
\(r\geq1\), \(q^{k-1}=q^{r-1}q^m\). The scaling remark and part~(2) of
the lemma give
\[
 \begin{aligned}
 \mathbb E\bigl(N(N-1)\cdots(N-r+1)\bigr)
 &=r!\,p\,q^{r-1}\sum_{m\geq0}\binom{m+r}{r}q^m\\
 &=r!\,p\,q^{r-1}p^{-(r+1)}=\frac{r!\,q^{r-1}}{p^r}.
 \end{aligned}
\]

(3) The orders \(r=1\) and \(r=2\) give \(\mathbb EN=1/p\) and
\(\mathbb E[N(N-1)]=2q/p^2=2(1-p)/p^2\). On \(S\),
\(x^2=x(x-1)+x\) with both summands nonnegative, so, exactly as in the
Poisson proof, \(\mathbb EN^2=2q/p^2+1/p<\infty\). Then
\cref{thm:p15-expectation-bounds} gives
\[
 \operatorname{Var}(N)=\frac{2q}{p^2}+\frac1p-\frac1{p^2}
 =\frac{2q+p-1}{p^2}=\frac{2q-q}{p^2}=\frac{1-p}{p^2}.
\]

(4) For \(k\geq1\),
\(\{N=k\}=A_1^c\cap\cdots\cap A_{k-1}^c\cap A_k\), and
\(\{N=\infty\}=\bigcap_{i\geq1}A_i^c\). These events are measurable.
For real \(a\), \(\{N\leq a\}\) is the finite union of the events
\(\{N=k\}\) with \(1\leq k\leq a\) (empty if \(a<1\)), so \(N\) is an
extended real random variable (\cref{def:p15-extended-random-variables}).
By part~(2) of \cref{thm:p15-independence-basic}, complementing the
first \(k-1\) of the independent events \(A_1,\ldots,A_k\) keeps them
independent, so
\(\mathbb P(N=k)=(1-p)^{k-1}p=pq^{k-1}\). The events
\(B_n=A_1^c\cap\cdots\cap A_n^c\) decrease to \(\{N=\infty\}\) and,
again by independence of complements, \(\mathbb P(B_n)=q^n\), which
tends to zero as in the proof of the lemma. Continuity from above
(\cref{thm:p15-measure-laws}) gives \(\mathbb P(N=\infty)=0\).
Put \(N'=N\) on \(\{N<\infty\}\) and \(N'=1\) on \(\{N=\infty\}\).
For real \(a\), \(\{N'\leq a\}\) is \(\{N\leq a\}\) when \(a<1\) and
\(\{N\leq a\}\cup\{N=\infty\}\) when \(a\geq1\), so \(N'\) is a
real random variable. It differs from \(N\) only on a null event, so
\(\mathbb P(N'=k)=\mathbb P(N=k)\) for every \(k\geq1\); thus \(N'\)
has the geometric law and parts~(2) and~(3) apply to it.

\emph{The case \(p=1\) stated after the theorem.} Then \(N=1\) a.s.
The descending product of order one is \(N=1\) a.s., with expectation
one. For \(r\geq2\) it contains the factor \(N-1=0\) a.s., so its
expectation is zero; this is \(r!\,0^{r-1}/1^r\) with \(0^0=1\).

\emph{The worked partial sums.} Let
\(s_K=\sum_{k=1}^Kk2^{-k}\). Then \(s_1=1/2=2-3\cdot2^{-1}\), and if
\(s_K=2-(K+2)2^{-K}\), then
\[
 \begin{aligned}
 s_{K+1}&=2-(K+2)2^{-K}+(K+1)2^{-K-1}\\
        &=2-\bigl(2K+4-K-1\bigr)2^{-K-1}=2-(K+3)2^{-(K+1)}.
 \end{aligned}
\]
At \(K=1,\ldots,5\) this gives \(1/2,1,11/8,13/8,57/32\).

\emph{The limits of \cref{def:p15-w02b-p15-001-limits}.}
For item~(1), \((1/2)(1/2-1)=-1/4\).

For item~(2), \(1/(k(k+1))=1/k-1/(k+1)\), so the partial sums
telescope to \(1-1/(K+1)\), which increase to one. By
\cref{thm:p15-countable-atomic} these masses define a probability on
the positive integers, and the identity map \(M\) has that law. By part~(1)
of the lemma, \(\mathbb EM=\sum_{k\geq1}1/(k+1)\). For \(m\geq1\),
split \(2\leq n\leq2^m\) into the blocks \(2^{j-1}<n\leq2^j\),
\(1\leq j\leq m\). Each block has \(2^{j-1}\) terms, each at least
\(2^{-j}\), so \(\sum_{n=2}^{2^m}1/n\geq m/2\). The partial sums are
unbounded and \(\mathbb EM=\infty\). Next,
\(\mathbb E[M(M-1)]=\sum_{k\geq1}(k-1)/(k+1)\), and for \(k\geq2\),
\((k-1)/(k+1)\geq1/3\) because \(3k-3\geq k+1\); so this sum is
infinite too. Since \(M^2\geq M\) on the positive integers,
\(\mathbb EM^2=\infty\), and \cref{def:p15-expectation-integral}
assigns a variance only when the second moment is finite. Subtraction
of infinities is not defined (\cref{def:p15-measure}).

For item~(3), at \(p=0\) every mass \(0\cdot1^{k-1}\) is zero, so the
masses sum to zero, not one. If every \(A_i\) is null, then
\(\mathbb P(N<\infty)=\mathbb P(\bigcup_iA_i)\leq\sum_i\mathbb P(A_i)=0\)
by countable subadditivity (\cref{thm:p15-measure-laws}), so
\(N=\infty\) almost surely.
\end{proof}
\end{annexitem}

\begin{annexitem}{annex:p15-w02b-p15-001-lp-results}{thm:p15-w02b-p15-001-lp-results}{The essential bound, the limit of norms and Minkowski's inequality}
\begin{proof}
Two facts about real powers are used. For fixed \(p>0\), the map
\(t\mapsto t^p\) is increasing on \([0,\infty)\), and so is its inverse
\(t\mapsto t^{1/p}\); both are continuous, and on \([0,\infty]\) they
are extended by \(\infty^p=\infty^{1/p}=\infty\). For fixed
\(c\in(0,\infty)\), \(c^{1/p}=\exp((\log c)/p)\to\exp(0)=1\) as
\(p\to\infty\), by continuity of the exponential function.

(1) Let \(C=\|f\|_\infty<\infty\). For each integer \(k\geq1\), the
definition of the infimum gives some \(M_k\geq0\) with
\(\mu(\{|f|>M_k\})=0\) and \(M_k<C+1/k\). Then
\(\{|f|>C+1/k\}\subseteq\{|f|>M_k\}\) is null. The union of these sets
over \(k\geq1\) is \(\{|f|>C\}\), which is null by countable
subadditivity (\cref{thm:p15-measure-laws}). So \(|f|\leq C\) a.e.,
and \(C\) itself belongs to the set whose infimum defines
\(\|f\|_\infty\).

Let \(M\geq0\). If \(\|f\|_\infty\leq M\), then \(\|f\|_\infty\) is
finite and \(|f|\leq\|f\|_\infty\leq M\) a.e. Conversely, if
\(|f|\leq M\) a.e., then \(\{|f|>M\}\) is contained in a null set and is
measurable, so it is null; thus \(M\) is admissible in the infimum and
\(\|f\|_\infty\leq M\).

For the product bound, suppose first \(\|g\|_\infty<\infty\). Then
\(|fg|=|f||g|\leq\|g\|_\infty|f|\) a.e., and part~(2) of
\cref{thm:p15-nonnegative-integral-laws} with nonnegative scaling gives
\(\int|fg|\,d\mu\leq\|g\|_\infty\int|f|\,d\mu=\|f\|_1\|g\|_\infty\).
Suppose \(\|g\|_\infty=\infty\). If \(\|f\|_1=0\), then \(f=0\) a.e. by
part~(1) of the same theorem, so \(fg=0\) a.e. and both sides are zero.
If \(\|f\|_1>0\), the right side is \(\infty\) and there is nothing to
prove.

(2) Let \(0<\mu(\Omega)<\infty\) and \(C=\|f\|_\infty\in[0,\infty]\).
For real \(p\geq1\), \(\|f\|_p\in[0,\infty]\).

Upper bound. If \(C<\infty\), then by part~(1) \(|f|^p\leq C^p\) a.e.,
so \(\int|f|^p\,d\mu\leq C^p\mu(\Omega)\) and
\(\|f\|_p\leq C\,\mu(\Omega)^{1/p}\). When \(\mu\) is a probability
measure this is \(\|f\|_p\leq C\), which is the last assertion of
part~(2); it holds trivially when \(C=\infty\). In general
\(C\mu(\Omega)^{1/p}\to C\), so \(\limsup_{p\to\infty}\|f\|_p\leq C\).
If \(C=\infty\) the same upper bound is trivial.

Lower bound. If \(C=0\) then \(\liminf_p\|f\|_p\geq0=C\). Otherwise let
\(M\) be real with \(0\leq M<C\). Then \(M\) is not admissible in the
infimum defining \(C\), so \(A=\{|f|>M\}\) has \(\mu(A)>0\), and
\(\mu(A)\leq\mu(\Omega)<\infty\). Since \(|f|^p\geq M^p\mathbf1_A\)
everywhere, \(\int|f|^p\,d\mu\geq M^p\mu(A)\) and
\(\|f\|_p\geq M\mu(A)^{1/p}\). As \(p\to\infty\) the right side tends
to \(M\), so \(\liminf_p\|f\|_p\geq M\). Letting \(M\) increase to
\(C\) (through all real values below \(C\) when \(C=\infty\)) gives
\(\liminf_p\|f\|_p\geq C\).

The two bounds give \(\lim_{p\to\infty}\|f\|_p=C\) in \([0,\infty]\).

(3) First \(1\leq p<\infty\). Pointwise
\(|f+g|\leq|f|+|g|\leq2\max(|f|,|g|)\), and since \(t\mapsto t^p\) is
increasing,
\[
 |f+g|^p\leq2^p\max(|f|^p,|g|^p)\leq2^p\bigl(|f|^p+|g|^p\bigr).
\]
If \(\|f\|_p=\infty\) or \(\|g\|_p=\infty\), the right side of
Minkowski's inequality is \(\infty\) and there is nothing to prove.
Assume both are finite. Integrating the pointwise bound with
nonnegative additivity and scaling gives
\(\int|f+g|^p\,d\mu\leq2^p(\|f\|_p^p+\|g\|_p^p)<\infty\), so
\(N\defeq\|f+g\|_p<\infty\). If \(N=0\) the inequality holds. If
\(p=1\), integrate \(|f+g|\leq|f|+|g|\) to get \(N\leq\|f\|_1+\|g\|_1\).

Let \(1<p<\infty\) and \(0<N<\infty\), and put \(q=p/(p-1)\), so
\(1/p+1/q=1\) and \((p-1)q=p\). Pointwise, with \(0^{p-1}=0\),
\[
 |f+g|^p=|f+g|\,|f+g|^{p-1}\leq|f|\,|f+g|^{p-1}+|g|\,|f+g|^{p-1}.
\]
By \cref{thm:p15-holder-integral} applied to \(f\) and the real function
\(|f+g|^{p-1}\),
\[
 \int|f|\,|f+g|^{p-1}\,d\mu\leq\|f\|_p
 \Bigl(\int|f+g|^{(p-1)q}\,d\mu\Bigr)^{1/q}
 =\|f\|_p\,N^{p/q}=\|f\|_p\,N^{p-1},
\]
and in the same way with \(g\) in place of \(f\). Integrating the
pointwise inequality with nonnegative additivity gives
\(N^p\leq(\|f\|_p+\|g\|_p)N^{p-1}\). Dividing by the positive finite
number \(N^{p-1}\) gives \(N\leq\|f\|_p+\|g\|_p\).

Now \(p=\infty\). If either size on the right is infinite there is
nothing to prove. Otherwise, by part~(1), \(|f|\leq\|f\|_\infty\)
outside a null set \(Z_1\) and \(|g|\leq\|g\|_\infty\) outside a null set
\(Z_2\). Outside the null set \(Z_1\cup Z_2\),
\(|f+g|\leq|f|+|g|\leq\|f\|_\infty+\|g\|_\infty\). By the
characterization in part~(1),
\(\|f+g\|_\infty\leq\|f\|_\infty+\|g\|_\infty\).

(4) Fix \(1\leq p\leq\infty\) and real \(c\).

Homogeneity. For finite \(p\),
\(\int|cf|^p\,d\mu=|c|^p\int|f|^p\,d\mu\) by nonnegative scaling, so
\(\|cf\|_p=|c|\,\|f\|_p\), with \(0\cdot\infty=0\) when \(c=0\). For
\(p=\infty\) and \(c\neq0\), \(\{|cf|>|c|M\}=\{|f|>M\}\) for every
\(M\geq0\), so the admissible bounds for \(cf\) are exactly \(|c|\)
times those for \(f\), and \(\|cf\|_\infty=|c|\,\|f\|_\infty\); for
\(c=0\) both sides are zero.

Vector space and seminorm. The zero function has size zero. If
\(f,g\in\mathcal L^p(\mu)\), then \(cf\in\mathcal L^p(\mu)\) by
homogeneity and \(f+g\in\mathcal L^p(\mu)\) by part~(3). Pointwise
operations on real functions satisfy the vector-space axioms, so
\(\mathcal L^p(\mu)\) is a real vector space, and homogeneity with
part~(3) makes \(\|\cdot\|_p\) a seminorm on it.

Zero size. For finite \(p\), \(\|f\|_p=0\) holds exactly when
\(\int|f|^p\,d\mu=0\), which by part~(1) of
\cref{thm:p15-nonnegative-integral-laws} holds exactly when
\(|f|^p=0\) a.e., that is, \(f=0\) a.e. For \(p=\infty\), part~(1)
with \(M=0\) shows that \(\|f\|_\infty=0\) exactly when \(|f|\leq0\)
a.e.

Equality a.e. Let \(f=g\) outside a null set \(Z\). For finite \(p\),
\(|f|^p=|g|^p\) a.e., so the integrals agree by part~(2) of
\cref{thm:p15-nonnegative-integral-laws}. For \(p=\infty\) and
\(M\geq0\), \(\{|f|>M\}\subseteq\{|g|>M\}\cup Z\) and conversely, so
one of these sets is null exactly when the other is; the admissible
bounds coincide.

The quotient. If \(f=f'\) and \(g=g'\) a.e., then \(f+g=f'+g'\) and
\(cf=cf'\) a.e., because the union of two null sets is null. So the
operations pass to classes, \(L^p(\mu)\) is a real vector space whose
zero is the class of functions equal to zero a.e., the size of a class
is well defined by the previous step, and it is a seminorm that
vanishes only on the zero class, that is, a norm.

\emph{The example \cref{def:p15-w02b-p15-001-lp-worked}.} Every subset
of \(\{1,2\}\) is an event. Only the empty set is null, so
\(\|f\|_\infty=\max(1,2)=2\), and
\(\|f\|_p^p=(1+2^p)/2\). Since \(2^p/2\leq(1+2^p)/2\leq2^p\), taking
\(p\)-th roots gives \(2\cdot2^{-1/p}\leq\|f\|_p\leq2\), and
\(2^{-1/p}\to1\). At \(p=1\), \((1+2)/2=3/2\); at \(p=2\),
\(((1+4)/2)^{1/2}=\sqrt{5/2}\), and \(1.58^2=2.4964<2.5<2.5281=1.59^2\).
Next \(\|g\|_p^p=\|h\|_p^p=1/2\), \(\|g\|_\infty=\|h\|_\infty=1\), and
\(g+h=1\) has size one for every \(p\). For \(1<p<\infty\),
\(1-1/p>0\) gives \(2^{1-1/p}>1\); at \(p=1\) it equals one.

\emph{The limits \cref{def:p15-w02b-p15-001-lp-limits}.} For item~(1),
\(\{1>M\}=\mathbb R\) has infinite measure for \(0\leq M<1\) and
\(\{1>1\}=\emptyset\), so \(\|1\|_\infty=1\), while
the integral of \(1^p=1\) with respect to Lebesgue measure is the
measure of \(\mathbb R\), which is infinite (\cref{thm:p15-lebesgue-unit-interval}). For item~(2), the same
computation as in the example gives \(\|g\|_p=\|h\|_p=2^{-1/p}\) and
\(\|g+h\|_p=1\); for \(0<p<1\), \(1-1/p<0\) and \(2^{1-1/p}<1\). For
item~(3), \(\{|f|>0\}=\{0\}\) has Lebesgue measure zero, so \(M=0\) is
admissible and \(\|f\|_\infty=0\), while \(f(0)=5\).
\end{proof}
\end{annexitem}

\begin{annexitem}{annex:p15-w02b-p15-001-approximations}{thm:p15-w02b-p15-001-approximations}{Dyadic sums, truncation, small sets and simple approximation}
\begin{proof}
(1) Each \(E_{n,m}=\{f\geq m/2^n\}\cap\{f<(m+1)/2^n\}\) is measurable.
For fixed \(n\), the sets \(E_{n,m}\), \(m\geq0\), are pairwise
disjoint, and since \(f\) is finite and nonnegative each \(x\) lies in
exactly one of them, namely the one with
\(m=\lfloor2^nf(x)\rfloor\). Hence
\[
 \varphi_n\defeq\sum_{m\geq0}\frac m{2^n}\mathbf1_{E_{n,m}}
 =2^{-n}\lfloor2^nf\rfloor
\]
pointwise. Termwise integration of a nonnegative series
(\cref{thm:p15-nonnegative-integral-laws}) and
\(\int c\mathbf1_E\,d\mu=c\mu(E)\) give
\(\int\varphi_n\,d\mu=\sum_{m\geq1}(m/2^n)\mu(E_{n,m})\); the term
\(m=0\) is zero. For real \(y\), \(2\lfloor y\rfloor\) is an integer at
most \(2y\), so \(\lfloor2y\rfloor\geq2\lfloor y\rfloor\). With
\(y=2^nf(x)\) this gives
\(\varphi_{n+1}=2^{-n-1}\lfloor2^{n+1}f\rfloor\geq
 2^{-n-1}\cdot2\lfloor2^nf\rfloor=\varphi_n\). Also
\(0\leq f-\varphi_n<2^{-n}\), so \(\varphi_n\uparrow f\) pointwise.
Monotone convergence (\cref{thm:p15-monotone-integral}) gives
\(\int\varphi_n\,d\mu\uparrow\int f\,d\mu\), which is the claim.

(2) For every \(x\), \(\min(f(x),n)\) is nondecreasing in \(n\) and
equals \(f(x)\) once \(n\geq f(x)\) when \(f(x)<\infty\), while it
equals \(n\to\infty=f(x)\) when \(f(x)=\infty\). Each \(f\wedge n\) is
measurable, since \(\{f\wedge n\leq a\}\) is \(\{f\leq a\}\) for
\(a<n\) and \(\Omega\) for \(a\geq n\). Monotone convergence gives the
claim.

(3) The function \(|g|\) is real, nonnegative and integrable. By
part~(2), \(\int|g|\wedge n\,d\mu\uparrow\int|g|\,d\mu<\infty\), so
there is an integer \(n\geq1\) with
\(\int|g|\,d\mu-\int|g|\wedge n\,d\mu<\epsilon/2\). Put
\(\delta=\epsilon/(2n)\) and let \(\mu(A)<\delta\). Write
\(|g|=|g|\wedge n+(|g|-|g|\wedge n)\), a sum of two nonnegative real
measurable functions. Nonnegative additivity gives
\(\int(|g|-|g|\wedge n)\,d\mu=\int|g|\,d\mu-\int|g|\wedge n\,d\mu\),
all terms being finite, and also
\[
 \int_A|g|\,d\mu=\int_A|g|\wedge n\,d\mu+\int_A(|g|-|g|\wedge n)\,d\mu.
\]
The first term is at most \(\int n\mathbf1_A\,d\mu=n\mu(A)<\epsilon/2\).
The second is at most \(\int(|g|-|g|\wedge n)\,d\mu<\epsilon/2\) by
monotonicity. Hence \(\int_A|g|\,d\mu<\epsilon\).

(4) Fix \(n\geq1\) and let
\(\psi_n(x)=\operatorname{sgn}(x)\min(n,2^{-n}\lfloor2^n|x|\rfloor)\)
for real \(x\), so \(\varphi_n=\psi_n\circ f\). The function \(\psi_n\)
takes its values in the finite set
\(V_n=\{j2^{-n}:j\text{ an integer},\ |j|\leq n2^n\}\), and each level
set \(\{\psi_n=c\}\) is a finite union of intervals: for
\(0<j<n2^n\) the level \(j2^{-n}\) is taken on \([j2^{-n},(j+1)2^{-n})\),
the level \(n\) on \([n,\infty)\), the level zero on
\((-2^{-n},2^{-n})\), and the negative levels symmetrically. Hence
\(\psi_n\) is Borel, \(\varphi_n\) is measurable as a composite
(\cref{thm:p15-measurable-operations}), and
\(\varphi_n=\sum_{c\in V_n}c\,\mathbf1_{\{\varphi_n=c\}}\) is a real
simple function. Since \(2^{-n}\lfloor2^n|f|\rfloor\leq|f|\),
\(|\varphi_n|\leq|f|\). The value \(\varphi_n(x)\) is nonzero exactly
when \(\lfloor2^n|f(x)|\rfloor\geq1\), that is, when
\(|f(x)|\geq2^{-n}\). On that set \(|f|^p\geq2^{-np}\), so
\(2^{-np}\mu(\{\varphi_n\neq0\})\leq\int|f|^p\,d\mu<\infty\).

Fix \(x\). When \(n\geq|f(x)|\), the minimum defining \(\varphi_n(x)\) is
\(2^{-n}\lfloor2^n|f(x)|\rfloor\), and
\(0\leq|f(x)|-|\varphi_n(x)|<2^{-n}\); as \(\varphi_n(x)\) and
\(f(x)\) have the same sign, \(|\varphi_n(x)-f(x)|<2^{-n}\). Thus
\(\varphi_n\to f\) pointwise, and \(h_n\defeq|\varphi_n-f|^p\to0\)
pointwise. Moreover
\(h_n\leq(|\varphi_n|+|f|)^p\leq2^p|f|^p\), an integrable function. By
part~(1) of \cref{thm:p15-dominated-bounded-convergence},
\(\int h_n\,d\mu=\|h_n-0\|_1\to0\), so
\(\|\varphi_n-f\|_p=(\int h_n\,d\mu)^{1/p}\to0\) by continuity of
\(t\mapsto t^{1/p}\) at zero.

\emph{The example \cref{def:p15-w02b-p15-001-approx-worked}.} For
\(f=U\) and \(1\leq m\leq2^n-1\), \(E_{n,m}=\{u\in(0,1):
m/2^n\leq u<(m+1)/2^n\}=[m/2^n,(m+1)/2^n)\), because
\(m/2^n>0\) and \((m+1)/2^n\leq1\). Its probability is its length
\(2^{-n}\) (\cref{thm:p15-lebesgue-unit-interval}). For \(m\geq2^n\),
\(m/2^n\geq1>u\), so \(E_{n,m}\) is empty. The sum is
\(4^{-n}\sum_{m=1}^{2^n-1}m=4^{-n}(2^n-1)2^n/2=(1-2^{-n})/2\) by
part~(1) of \cref{thm:p15-w02b-p15-001-uniform-moments}. At
\(n=1,2,3\) this is \(1/4\), \(3/8\) and \(7/16\).

\emph{The limits \cref{def:p15-w02b-p15-001-approx-limits}.} For
item~(1), \(f(1)=\infty\) satisfies no inequality \(f<(m+1)/2^n\), so
every \(E_{n,m}\) is empty; \(\int f\,d\mu\geq\int n\mathbf1_\Omega\,d\mu=n\)
for every \(n\), so \(\int f\,d\mu=\infty\); and \(f\wedge n=n\) on
\(\Omega\) has integral \(n\). For item~(2), \(\mu\) is defined by
\cref{thm:p15-countable-atomic} from the masses \(2^{-k}\), and
\(\int_{\{k\}}|g|\,d\mu=2^k\cdot2^{-k}=1\). Given \(\delta>0\), choose
\(k\) with \(2^{-k}<\delta\); the set \(A=\{k\}\) has
\(\mu(A)<\delta\) and \(\int_A|g|\,d\mu=1\), which is not below
\(\epsilon=1\). Termwise integration gives
\(\int|g|\,d\mu=\sum_k1=\infty\). For item~(3), let \(\varphi\) be real
simple whose nonzero set has finite Lebesgue measure. Then
\(\{\varphi=0\}\) has infinite measure, since otherwise \(\mathbb R\),
the union of the two sets, would have finite measure. On it \(|\varphi-f|=1\), so
for \(0\leq M<1\) the set \(\{|\varphi-f|>M\}\) is not null, and
\(\|\varphi-f\|_\infty\geq1\).
\end{proof}
\end{annexitem}
